\documentclass[12pt, a4paper]{article}
\usepackage[utf8]{inputenc}
\usepackage[T1]{fontenc}
\usepackage{amsmath, amssymb, amsthm}
\usepackage{geometry}
\usepackage{setspace}
\usepackage{hyperref}

\geometry{margin=1in}
\setstretch{1.15}

\title{Fixed-Point Semantics in Discrete Temporal Logic: A Constructivist Ontology of the Mind}
\author{Independent Researcher}
\date{\today}

\newtheorem{formalism}{Formalism}

\begin{document}

\maketitle

\begin{abstract}
This paper presents a formal ontology for a discrete temporal logic engine governed by the ``Indexical Present'' ($I_k$). Unlike classical models of universal reality, this system defines truth as a localized, observer-dependent manifestation of information called ``events.'' Addressing the paradoxes inherent in continuous models of time, we propose a system utilizing a monotonic Positive Logic to ensure that every discrete time step—here defined as a \textit{Chronon}—reaches a stable Least Fixed Point via the Knaster-Tarski theorem. The ontology bifurcates the universe into two domains: the ``Noumenal'' (The Raw; the inconceivable Absolute) and the ``Phenomenal'' (The Abstract; the thinker's conceptual model). We argue that the thinker is a ``Real Void''—a real activity of being producing unreal, information-based content. This framework offers a rigorous resolution to the ``Paradox of the Abstract,'' positioning scientific consistency not as a map of the Absolute, but as the coherence of a model maintained against the friction of the Real.
\end{abstract}

\section{Events and the Indexical Present}

We define the \textit{Chronon} as the abstraction of a discrete temporal step. It is a unit of simultaneity that serves as the interval wherein all logical implications reach a \textit{fixed point}. Crucially, the Chronon is not a duration measured \textit{within} physical time; rather, it is the atomic logical operation from which physical time \textit{emerges}. It possesses ``logical depth'' (computational time) but zero ``physical duration.''

\begin{formalism}
Let $\mathbb{T} = \mathbb{N}_0$ be the set of discrete time steps (Chronons).
Let $k \in \mathbb{T}$ denote the current temporal index.
\end{formalism}

An \textit{event} is an atomic unit of information existing in the \textbf{Indexical Present} ($I_k$)—the logical equivalent of the ``Observed Now.'' Events are the language of knowledge; they are separable, conceivable units of information. However, they do not constitute ``The Raw'' (the Absolute), which remains inconceivable.

An event is either \textit{caused} or \textit{implied}. While logical implications manifest instantaneously within the Chronon, \textbf{Causation} is the bridge between the Noumenal and the Phenomenal. History provides the \textit{context} for an event, but the \textit{precipitating force} is the Absolute.

\begin{formalism}
Let $\mathcal{E}$ be the finite set of all possible event symbols.
The \textbf{Indexical Present} at step $k$ is a set $I_k \subseteq \mathcal{E}$.
Let $\mathfrak{I}: \mathcal{P}(\mathcal{E}) \to \mathcal{P}(\mathcal{E})$ be the \textbf{Implication Function} (deduction within $k$).
\end{formalism}

The Indexical Present is singular and persistent; it is the vessel that holds the \textit{existence} of events. Events, together with their implications, inhabit $I_k$ for exactly one Chronon—the logical instant required for the set to saturate and stabilize. The set of stable events is then committed to \textit{History}.

\begin{formalism}
Let $\mathcal{E}$ be the Universe of Events and $k \in \mathbb{N}_0$ be the current index.
\textbf{History} $H_k$ is defined as the immutable, ordered tuple of event sets up to Chronon $k-1$.

$$H_k := \langle I_0, I_1, \dots, I_{k-1} \rangle \in (\mathcal{P}(\mathcal{E}))^{k}$$

\textbf{The Crystallization (Transition):}
$$H_{k+1} = H_k \frown \langle I_{k} \rangle$$

\textbf{The Primary Mover ($k=0$):}
At $k=0$, $H_0 = \langle \rangle$. The first Chronon $I_0$ relies on an initial injection of events from the Noumenal source (The Raw), setting the phases of the Chronons in motion.
\end{formalism}

\section{Constructive Logic and Implication}

Implications within the Indexical Present operate under a strictly \textbf{Constructive Positive Logic}. Existence is strictly additive; truth is established by construction (presence), not by the negation of falsehood. There is no operation of 'not event' within the active Chronon. Negation is restricted solely to queries against the immutable \textit{history} (classical Boolean validation).

\begin{formalism}
The Implication Function $\mathfrak{I}(S)$ is \textbf{Monotonic}:
$$\forall A, B \subseteq \mathcal{E} : A \subseteq B \implies \mathfrak{I}(A) \subseteq \mathfrak{I}(B)$$
Negation ($\neg e$) is undefined for $e \in I_k$.
Negation is valid only for predicates $P(H)$ over history, where the Law of Excluded Middle holds because $H$ is a completed object.
\end{formalism}

In this system, \textit{contradictions are structurally impossible} within the present moment. The absence of an event is not a logical value that conflicts with presence; it is simply silence.

\subsection{The Dual Logic Frame}

To maintain the integrity of the Indexical Present while allowing for complex reasoning, we distinguish between two logical domains:

\begin{enumerate}
    \item \textbf{The Domain of Existence (The Consequent):}
    \begin{itemize}
        \item \textbf{Target:} The Indexical Present ($I_k$).
        \item \textbf{Logic:} Strictly Positive/Constructive. Events are atomic tokens.
    \end{itemize}

    \item \textbf{The Domain of Validation (The Antecedent):}
    \begin{itemize}
        \item \textbf{Source:} History ($H_k$) and State ($\mathcal{A}_k$).
        \item \textbf{Logic:} General Boolean Logic with Temporal Operators. Because history is immutable, queries against it return definitive \textbf{True/False} values.
    \end{itemize}
\end{enumerate}

The rule $\psi \vdash e$ functions as a transducer: it converts a Truth Value (from Validation) into an Event (in Existence).

\section{Stability and the Least Fixed Point}

This system eliminates logical instability (oscillation) through Idempotence and Monotonicity. An implication is a constructive force. To guarantee stability, we define the Chronon calculation as an iterative process.

\begin{formalism}
\textbf{1. The Injection Step ($S_0$):}
The initial state of the Chronon is populated by events precipitated by the Absolute, conditioned by History ($H_k$) and Agency ($A$).
$$S_0 = \text{Absolute}(H_k, \Delta \mathcal{A}_k)$$

\textbf{2. The Saturation Step ($S_{i} \to S_{i+1}$):}
We apply the Implication Function $\mathfrak{I}(S)$ iteratively.
Let $\Phi(S) = S \cup \mathfrak{I}(S)$ be the Step Function.
$$S_{i+1} = \Phi(S_i)$$

\textbf{3. Convergence:}
Since $S \subseteq \Phi(S)$ (Inflationary) and $\Phi$ is Monotonic, by the \textbf{Knaster-Tarski Theorem}, the sequence reaches a \textbf{Least Fixed Point}:
$$I_k = \text{lfp}(\Phi, S_0) = \bigcup_{i=0}^{\infty} \Phi^i(S_0)$$
\end{formalism}

\section{Ontic Friction: The Role of Absurdity}

The co-existence of implied events in the Indexical Present guarantees logical stability (termination), but it does not guarantee semantic consistency.

We propose that Absurdity is not merely a system error, but the \textbf{Detection of Reality}. Following the premise that the Abstract Model is a construct of the Thinker, \textbf{Consistency} is the only metric available to measure the quality of the model. When the Absolute precipitates events that contradict the logic of the Abstract Model (e.g., \texttt{DOOR\_CLOSED} $\land$ \texttt{DOOR\_OPEN}), the system experiences \textbf{Ontic Friction}.

\begin{formalism}
Let $\mathcal{K}$ be a set of Consistency Constraints.
A Chronon $I_k$ is \textbf{Valid} iff: $I_k = \text{lfp}(\Phi)$.
A Chronon $I_k$ is \textbf{Absurd} iff: $\exists \kappa \in \mathcal{K} : I_k \not\models \kappa$.
\end{formalism}

Absurdity is the friction between the \textbf{Abstract Expectation} and the \textbf{Noumenal Manifestation}. It serves as the primary sensor for the ``Real.''

\section{State: The Lens of History}

We distinguish between the \textit{Object} of history and the \textit{Interpretation} of history. As $k \to \infty$, the volume of History becomes intractable. The observer requires \textbf{State}—a derived \textbf{Interpretation} of History. This acts as a compression algorithm, filtering temporal depth into a manageable set of persistent properties.

\begin{formalism}
Let $\mathbb{S}$ be the set of possible States.
Let $\Psi$ be the \textbf{Interpretation Function} (The Lens).
$$S_k = \Psi(H_k)$$
\end{formalism}

The Indexical Present is the source of perceived Truth (Events). History is the sustainer of Meaning (State).

\section{Relational Quantum Mechanics and the ``I Am''}

This ontology aligns with \textbf{Relational Quantum Mechanics}, positing that there are no observer-independent states. The Indexical Present ($I_k$) is not a ``God's Eye View'' but a local state relative to the observer's frame.

However, there is one absolute truth that is perceived without the necessity of an event: \textbf{The truth of one’s own existence.} For the observer, this truth is immediate and absolute. All other truths are accessible only as knowledge (events) within the Indexical Present.

\begin{formalism}
Let $\mathcal{O}$ be the set of Observers. Truth is strictly indexed by the observer:
$$Truth(e, o, k) \iff e \in I_{k, o}$$

\textbf{Ternary Logic of the Chronon:}
Let $V(p)$ interpret the fixed point $I_k$:
$$V(p) = \begin{cases} \text{True} & \text{if } p \in I_{k, o} \\ \text{False} & \text{if } \exists e \in I_{k, o} : \{e, p\} \text{ violates } \mathcal{K} \\ \textbf{Unknown} & \text{otherwise} \end{cases}$$
\end{formalism}

\textbf{The Nature of Unknown:} The value \textbf{Unknown} is not merely a lack of data; it is the formal placeholder for the \textbf{Noumenal}. It represents the ``Raw'' reality that exists but has not manifested as a conceivable event.

In the case of Schrödinger’s Cat:
\begin{enumerate}
    \item \textbf{In the Cat's Frame:} The event triggers immediately. The cat is an observer with its own $I_{k, cat}$ and its own certainty of existence.
    \item \textbf{In the Scientist's Frame:} No event \texttt{CAT\_DEAD} or \texttt{CAT\_ALIVE} has manifested. The value is formally \textbf{Unknown}—the Scientist faces the Noumenal Latency.
\end{enumerate}

\section{The Abstract and the Real Void}

The \textbf{Abstract Model} is the observer’s internal representation of the world. Within this domain, the observer acts as a Thinker. We posit a fundamental ontological distinction:

\begin{enumerate}
    \item \textbf{The Activity is Real:} The act of conceiving possesses the quality of Absolute Being.
    \item \textbf{The Content is Unreal:} The content produced (categories, dimensions, objects) remains purely conceptual.
\end{enumerate}

Thus, the Thinker is a \textbf{Real Void}: a real activity of being producing unreal, information-based content.

\subsection{Agency as Transcendental Correspondence}
We distinguish between events derived from deduction and events that enter the system via \textbf{Agency}. Agency is not a mechanical input but a \textbf{Singularity}.

The Abstract (the Mind) possesses structure but no force. The Absolute (the Raw) possesses Being but no structure. The transition from the Thinker's \textit{Intent} (a change in the Abstract) to a \textit{Realized Event} represents a **Transcendental Correspondence**: the Absolute reflects the state change of the Abstract into the Indexical Present without a describable, conceivable intermediate linkage.

\begin{formalism}
Let $\mathcal{A}$ be the Abstract State Space.
$$e \in I_{k+1} \iff \text{Absolute}(\Delta \mathcal{A}_k) \to e$$
\end{formalism}

We do not ``understand'' Agency; we rely on it. We are dependent on the Absolute to produce events that allow us to maintain the consistency of our Abstract Model.

\section{Three Phases of the Cycle}

The engine operates in a rhythmic cycle ensuring logical simultaneity within a Chronon while enforcing succession across them.

\begin{enumerate}
    \item \textbf{Phase I: Deduction (Simultaneity):} Events precipitated by the Absolute (conditioned by History or Agency) manifest. The Implication Function expands the set to the Least Fixed Point.
    \item \textbf{Phase II: Crystallization (Persistence):} The stable set $I_k$ is appended to History ($H_{k+1}$). The Interpretation Function updates the State ($\mathcal{A}_{k+1}$).
    \item \textbf{Phase III: Manifestation (Succession):} The Thinker modifies the Abstract State (Intent). The Absolute translates this into new events for the next Chronon ($C_{k+1}$).
\end{enumerate}

\section{Time and Dimension as Constructs}

Time and space are not properties of the Noumena (The Raw), but are internal constructs of the Phenomenal Model. Dimensions are properties of the Interpretation Function $\Psi$. Since the Abstract Model is a conceptual space, its dimensionality is restricted to the Thinker’s capacity to categorize history. The Noumenal remains non-dimensional.

\section{The Paradox of the Abstract}

The Thinker operates within an Abstract Model based on interpreted History. The only way to distinguish between events originating from the Absolute (Reality) and events implied by the Abstract (Imagination) is \textbf{Consistency}.

As the Thinker refines their internal implications, the model becomes increasingly consistent. A high consistency metric ($\Gamma$) poses the risk of \textbf{Ontological Blindness}—mistaking the map for the territory.

\begin{formalism}
Let $E_{abs}$ be events precipitated by the Absolute and $E_{inf}$ be events implied by the Abstract.
\textbf{Consistency Metric ($\Gamma$):}
$$\Gamma = \text{Consistent}(E_{abs} \cup E_{inf})$$
\end{formalism}

If $\Gamma$ is perfectly high, the observer cannot distinguish the Model from the Noumena. Only through the friction of \textbf{Absurdity} ($\Gamma < 1$) is the ``Real'' detected.

\section{Conclusion}

This paper formalizes a discrete temporal ontology that shifts focus from ``Objective Reality'' to ``Local Consistency.'' By defining the Chronon as a logical fixed-point calculation and the Indexical Present as the exclusive vessel of constructive truth, we acknowledge that an observer can never grasp ``The Raw'' (the Absolute). We are, instead, inhabitants of the Abstract—a conceptual scaffolding built from the persistent crystallization of ephemeral events into history.

The rigor of this system lies in its use of Positive Logic. By barring negation from the Now and allowing for the manifestation of Absurdity, the engine does not seek to prevent paradox, but to reveal it. Inconsistency becomes the primary sensor for the ``Real''; it is the friction that reminds the Thinker that their Abstract model is not the Absolute.

Ultimately, this ontology defines the Thinker not as a master of the universe, but as a \textbf{Steward of a Model}. Agency—the transformation of intent into action—is the point of deepest mystery, where the Absolute reflects the changes of our unreal imagination back into the real manifestations of our experience.

\end{document}

